\scanpage{158}
% archive foliation: 157
\uline{Catégorie galoisienne}

(i) Une catégorie $\mathcal C$

(ii) Un foncteur $F:\mathcal C\to\text{Ensembles finis}$

Axiomes:

\margin{\unsure{Poser} plutôt en axiome l'\uline{existence} d'un $F$ ayant les propriétés: \deleted{\unsure{Soit}} $F$ permute aux sommes, aux produits, aux \unsure{quotients par} groupes finis, et est tel qu'une \unsure{injection} $F(X)\to F(Y)$ \unsure{($X$ connex.)} $\Rightarrow X\overset{\sim}{\longrightarrow}Y$.}

(Gal 1) Dans $\mathcal C$, les sommes et les produits \deleted{de deux objets}\inserted{finis} existent [d'où l'existence de $\underline0$ et $\underline1$]

\deleted{(Gal 3) $X\in\mathcal C$, $F(X)=\varnothing\Longrightarrow X$ est une unité à gauche (soit $\underline0$)}

(Gal 4) $F$ permute à la somme et au produit

$X\in\mathcal C$ est dit \deleted{\uline{irréductible}}\inserted{connexe} si $X=X'\amalg X''$ implique $X'\simeq\underline0$ ou $X''\simeq\underline0$

(Gal 2) $X\in\mathcal C$ \deleted{$X\ne\underline0$} implique
\[
X\simeq\coprod_{i\in I}X_i,\quad \inserted{X_i\ne\underline0},
\]
et tte décomposition $X\simeq X'\amalg X''$ s'explique
\[
X'\simeq\coprod_{i\in I'}X_i,\quad X''=\coprod_{i\in I''}X_i,\quad
I'\cap I''=\varnothing,\quad I'\cup I''=I.
\]
Les $X_i$, déterminés de façon essentiellement unique [avec les $X_i\to X$] sont appelés les \uline{comp.} \deleted{\uline{irréductibles}}\inserted{connexes} de $X'$ [N.B. $X=\underline0\Longleftrightarrow I=\varnothing$]

(\uuline{Gal 3}) Soit $f:X\to Y$ avec $Y$ \deleted{irréductible}\inserted{connexe $\ne\underline0$}.
Alors les composantes \inserted{connexes} $X_i$ de $X$ telles que $f$ induise un \uline{isom.} $f_i:X_i\overset{\sim}{\longrightarrow}Y$ sont en correspondance biunivoque avec les \uline{sections} $g_i:Y\to X$ \deleted{$\unsure{f}$}
[hom. tels que $fg=\operatorname{id}_Y$], $g_i$ étant le composé \deleted{\unsure{de}}
\diagram{\begin{tikzcd}
Y \arrow[r,"f_i^{-1}"] & X_i \arrow[r,"\text{inj. can.}"] & X
\end{tikzcd}}
[en d'autres termes, toute section se factorise]
[N.B. il suffit de le vérifier si $X=Z\times Y$, $f=p_2$]

\margin{Gal 3 et Gal 5 ensemble sont impliqués par:
Si $f:X\to Y$ est tel que $F(X)\to F(Y)$ soit \uline{injectif} (\deleted{\unsure{$Y$ connexe}}\inserted{$X,\unsure{Y}$ connexes}) alors c'est un isomorphisme \inserted{de} $X$ sur une composante connexe $Y_i$.}

(Gal 5) Dans les conditions de Gal 3, si \deleted{\unsure{le morphisme}}\inserted{$X,Y$ connexes et} $f:X\to Y$ induit $F(X)\overset{\sim}{\longrightarrow}F(Y)$, c'est un isomorphisme
[implique $F(X)=\varnothing\longrightarrow X=\underline0$]
[Il suffit de vérifier Gal 5 quand $X$ est une composante connexe d'un produit $Z\times Y$, $f$ étant induit par $p_2$]

(\uuline{Gal 6}) Soit $X\in\mathcal C$, $X$ connexe $\ne\underline0$, et \uline{galoisien}, \deleted{\unsure{c.à.d.}} i.e. \deleted{\unsure{Alors}} $G=\operatorname{Aut}(X)$ opère \uline{transitivement} dans $F(X)$.
\deleted{Tte composante \unsure{connexe d'un revêtement est galoisienne}}
Alors pour tout sous-groupe $H$ de $G$, $X/H$ dans $\mathcal C$ \uline{existe}, et
\[
\unsure{F(X/H)|E(X/H)}\simeq F(X)/H.
\]
